\subsection{定理1的证明：第二部分}

再次考虑正合列

\[
0 \longrightarrow t ^ {r} / t ^ {r + 1} \xrightarrow {i _ {r}} A _ {0} / t ^ {r + 1} \xrightarrow {p _ {r}} A _ {0} / t ^ {r} \longrightarrow 0\tag{16}
\]

以及相伴的挠模正合列

\[
\begin{array}{l} \text {Tor} _ {1} ^ {\mathrm{A} _ {0}} (\mathrm{A} _ {0} / t ^ {r + 1}, \mathrm{M}) \longrightarrow \mathrm{Tor} _ {1} ^ {\mathrm{A} _ {0}} (\mathrm{A} _ {0} / t ^ {r}, \mathrm{M}) \\ \longrightarrow (t ^ {r} / t ^ {r + 1}) \otimes_ {\mathrm{A} _ {0}} \mathrm{M} \xrightarrow {i , \otimes 1 _ {\mathrm{M}}} (\mathrm{A} _ {0} / t ^ {r + 1}) \otimes_ {\mathrm{A} _ {0}} \mathrm{M} \xrightarrow {p , \otimes 1 _ {\mathrm{M}}} (\mathrm{A} _ {0} / t ^ {r}) \otimes_ {\mathrm{A} _ {0}} \mathrm{M} \longrightarrow 0. \end{array} \tag {18}
\]

\(p_r \otimes 1_M\) 的核与 \(x^r M/x^{r+1} M\) 等同；此外， \(t^r/t^{r+1}\) 被 \(T_i\) 零化，并与次数为 \(r\) 的 \(A_0\) 的齐次分量等同，因此同态 \((t^r/t^{r+1}) \otimes_{A_0} M \to x^r M/x^{r+1} M\) （它由 \(i_r \otimes 1_M\) 导出）在次数 \(r\) 处与同态 \(\beta_M^x\) 的齐次分量等同。因此由正合列（18）可知，条件（v）等价于

（v'）：同态 \(\operatorname{Tor}_{1}^{\mathbf{A}_{0}}(p_{r},1_{\mathbf{M}})\) : \(\operatorname{Tor}_{1}^{\mathbf{A}_{0}}(\mathbf{A}_{0}/t^{r+1},\mathbf{M})\to\operatorname{Tor}_{1}^{\mathbf{A}_{0}}(\mathbf{A}_{0}/t^{r},\mathbf{M})\) 对所有 \(r\geqslant1\) 是满射的。

我们仍需证明蕴含关系 (v) \(\Rightarrow\) (i) 在诸模

\[
\mathbf {M} / (x _ {1} \mathbf {M} + \dots + x _ {i - 1} \mathbf {M})
\]

对 \(\mathfrak{x}\) -进拓扑是分离的时成立（ \(1 \leqslant i \leqslant n\) ）。设 \(\overline{\mathbf{M}}\) 为 A-模 \(\mathbf{M}/x_1 \mathbf{M}\) 。根据定义，序列 \(\mathbf{x}\) 对 \(\mathbf{M}\) 是正则的当且仅当 \((x_1)_\mathbf{M}\) 是单射且序列 \(\mathbf{x}' = (x_2, \ldots, x_n)\) 对 \(\overline{\mathbf{M}}\) 是正则的。通过对 \(n\) 作归纳推理，因此只需证明：若 \(\mathbf{M}\) 对 \(\mathfrak{x}\) -进拓扑是分离的且 \(\beta_{\mathbf{M}}^{\mathbf{x}}\) 是双射，则 \((x_1)_\mathbf{M}\) 是单射且 \(\beta_{\mathbf{M}}^{\mathbf{x}'}\) 是双射。现在， \(\beta_{\mathbf{M}}^{\mathbf{x}}\) 的双射性特别地蕴含了以 \(x_1\) 为比率、到 \(\bigoplus_{r} x^r \mathbf{M}/x^{r+1} \mathbf{M}\) 上的位似是单射的，从而 \(\text{Ker } (x_1)_\mathbf{M} \subset \bigcap_{i} x^i \mathbf{M}\) ，因此 \((x_1)_\mathbf{M}\) 是单射，如果 \(\mathfrak{x}\) -进拓扑在 \(\mathbf{M}\) 上是分离的。

因此我们归结为证明：如果 \((x_{1})_{\mathbf{M}}\) 是单射且 M 满足条件 \((\mathbf{v}')\) ，则 \(\overline{M}\) 满足条件 \((\mathbf{v}')\) （关于序列 \(x'\) ）。

根据假设，我们有一个正合列

\[
0 \longrightarrow \mathrm{M} \xrightarrow {(x _ {1}) _ {M}} \mathrm{M} \longrightarrow \overline {{\mathrm{M}}} \longrightarrow 0;
\]

令 \(\overline{\mathbf{A}}_{0} = \mathbf{A}_{0}/\mathbf{T}_{1}\mathbf{A}_{0},\quad\overline{\mathbf{t}} = \mathbf{t}\overline{\mathbf{A}}_{0}\) 。设 \(q:L\to M\) 为 \(\mathbf{A}_{0}\) -模 M 的一个自由分解；由于比率为 \(\mathbf{T}_{1}\) 的位似在 \(\mathbf{A}_{0}\) 中是单射，它在 L 中也是单射，并且我们有一个复形的正合列

\[
0 \longrightarrow \mathrm{L} \xrightarrow {(x _ {1}) _ {\mathrm{L}}} \mathrm{L} \longrightarrow \overline {{\mathrm{L}}} \longrightarrow 0
\]

其中 \(\overline{\mathbf{L}} = \mathbf{L} / x_1\mathbf{L}\) ，以及一个交换图

\[
\begin{array}{c} 0 \longrightarrow \mathrm{L} \stackrel {(x _ {1}) _ {\mathrm{L}}} {\longrightarrow} \mathrm{L} \longrightarrow \overline {{\mathrm{L}}} \longrightarrow 0 \\ \Big \downarrow^ {q} \quad \Big \downarrow^ {q} \quad \Big \downarrow^ {\overline {{q}}} \\ 0 \longrightarrow \mathrm{M} \stackrel {(x _ {1}) _ {\mathrm{M}}} {\longrightarrow} \mathrm{M} \longrightarrow \overline {{\mathrm{M}}} \longrightarrow 0. \end{array}
\]

由于 \(q\) 是一个拟同构， \(\overline{q}:\overline{\mathbf{L}}\to\overline{\mathbf{M}}\) 是 \(\overline{\mathbf{A}}_{0}\) -模 \(\overline{\mathbf{M}}\) 的一个自由分解（第 30 页，推论 1）。对每个 \(\overline{\mathbf{A}}_{0}\) -模 P，存在一个典范同构

\[
\mathbf {P} \otimes_ {\mathbf {A} _ {0}} \mathbf {L} \rightarrow \mathbf {P} \otimes_ {\overline {{{\mathbf {A}}}} _ {0}} \overline {{{\mathbf {L}}}},
\]

由此，通过取同调，得到一个同构

\[
\varphi_ {P}: \operatorname{Tor} _ {1} ^ {\mathbf {A} _ {0}} (\mathbf {P}, \mathbf {M}) \rightarrow \operatorname{Tor} _ {1} ^ {\overline {{\mathbf {A}}} _ {0}} (\mathbf {P}, \overline {{\mathbf {M}}}).
\]

若 \(u: \mathbf{P} \to \mathbf{P}'\) 是一个 \(\overline{\mathbf{A}}\) -同态，则

\[
\varphi_ {\mathrm{P} ^ {\prime}} \circ \operatorname{Tor} _ {1} ^ {\mathbf {A} _ {0}} (u, 1 _ {\mathrm{M}}) = \operatorname{Tor} _ {1} ^ {\overline {{\mathbf {A}}} _ {0}} (u, 1 _ {\mathrm{M}}) \circ \varphi_ {\mathrm{P}}.
\]

既然如此，假设条件 (v') 对 M 成立，下面我们来证明它对 \(\overline{M}\) 成立。设 \(r\) 为一个整数 \(\geqslant 1\) ；我们有一个行正合的交换图

\[
\begin{array}{c} 0 \longrightarrow \mathrm{A} _ {0} / t ^ {r} \xrightarrow {\mathrm{T} _ {1}} \mathrm{A} _ {0} / t ^ {r + 1} \longrightarrow \overline {{\mathrm{A}}} _ {0} / \overline {{t}} ^ {r + 1} \longrightarrow 0 \\ \downarrow^ {p _ {r - 1}} \quad \downarrow^ {p _ {r}} \quad \downarrow^ {\bar {p} _ {r}} \\ 0 \longrightarrow \mathrm{A} _ {0} / t ^ {r - 1} \xrightarrow {\mathrm{T} _ {1}} \mathrm{A} _ {0} / t ^ {r} \longrightarrow \overline {{\mathrm{A}}} _ {0} / \overline {{t}} ^ {r} \longrightarrow 0 \end{array}
\]

由此我们推出一个具有正合行的交换图

\[
\begin{array}{c c c c c} \operatorname{Tor} _ {1} ^ {\mathbf {A} _ {0}} (\mathbf {A} _ {0} / t ^ {r + 1}, \mathbf {M}) & \longrightarrow & \operatorname{Tor} _ {1} ^ {\mathbf {A} _ {0}} (\overline {{\mathbf {A}}} _ {0} / \overline {{t}} ^ {r + 1}, \mathbf {M}) & \longrightarrow & \mathbf {M} / x ^ {r}   \mathbf {M} \xrightarrow {x _ {1}} \mathbf {M} / x ^ {r + 1}   \mathbf {M} \\ \operatorname{Tor} _ {1} (p _ {r}, 1) & \downarrow & \operatorname{Tor} _ {1} (\bar {p} _ {r}, 1) & \downarrow & \downarrow \\ \operatorname{Tor} _ {1} ^ {\mathbf {A} _ {0}} (\mathbf {A} _ {0} / t ^ {r}, \mathbf {M}) & \longrightarrow & \operatorname{Tor} _ {1} ^ {\mathbf {A} _ {0}} (\overline {{\mathbf {A}}} _ {0} / \overline {{t}} ^ {r}, \mathbf {M}) & \longrightarrow & \mathbf {M} / x ^ {r - 1}   \mathbf {M} \xrightarrow {x _ {1}} \mathbf {M} / x ^ {r}   \mathbf {M}. \end{array}
\]

现在乘以 \(x_{1}\) 定义了从 \(\mathbf{M}/\mathfrak{x}^{r}\mathbf{M}\) 到 \(\mathbf{M}/\mathfrak{x}^{r+1}\mathbf{M}\) 的一个单射：通过对 \(r\) 归纳，这由正合列

\[
0 \rightarrow x ^ {r - 1} M / x ^ {r} M \rightarrow M / x ^ {r} M \rightarrow M / x ^ {r - 1} M \rightarrow 0
\]

以及比率为 \(x_{1}\) 的位似到 \(\bigoplus_{r\geqslant 0}(\mathfrak{x}^{r}\mathrm{M}/\mathfrak{x}^{r+1}\mathrm{M})\) 上的单射性立即得出。因此条件 (v) 蕴含同态 \(\operatorname{Tor}_{1}^{\mathbf{A}_{0}}(\overline{p}_{r},1_{\mathbf{M}})\) 对所有 \(r\geqslant 1\) 是满射的。通过与同构 \((\varphi_{\overline{\mathbf{A}}_{0}/\overline{\mathbf{t}}^{r+1}})^{-1}\) 和 \(\varphi_{\overline{\mathbf{A}}_{0}/\overline{\mathbf{t}}^{r}}\) 复合，我们推出同态

\[
\operatorname{Tor} _ {1} ^ {\overline {{\mathbf {A}}} _ {0}} \left(\bar {p} _ {r}, 1 _ {\mathbf {M}}\right): \operatorname{Tor} _ {1} ^ {\overline {{\mathbf {A}}} _ {0}} \left(\overline {{\mathbf {A}}} _ {0} / \bar {t} ^ {r + 1}, \mathbf {M}\right)\rightarrow \operatorname{Tor} _ {1} ^ {\overline {{\mathbf {A}}} _ {0}} \left(\overline {{\mathbf {A}}} _ {0} / \bar {t} ^ {r}, \mathbf {M}\right)
\]

对 \(r \geqslant 1\) 是满射的，从而得到 \(\overline{\mathbf{M}}\) 的条件 (v')。这就完成了定理的证明。

